\clearpage
\thispagestyle{empty}

\begin{center}
    {\Large \textbf{Zusammenfassung}}
\end{center}

\vspace{0.8cm}

\noindent Die automatische Erkennung potenzieller Brutstätten von \textit{Aedes aegypti} in urbanen Gebieten ist ein wichtiger Bestandteil datengetriebener Ansätze zur Unterstützung der Vektorkontrolle. UAV-basierte Bild- und Videodaten ermöglichen eine effiziente Erfassung größerer städtischer Bereiche, ihre automatische Analyse ist jedoch anspruchsvoll, da relevante Objekte häufig klein, teilweise verdeckt und in komplexe Hintergrundstrukturen eingebettet sind. Zusätzlich erfordert das Training leistungsfähiger Objekterkennungsmodelle große Mengen hochwertiger Annotationen, deren manuelle Erstellung bei Videodaten sehr zeitaufwendig ist.

Diese Arbeit untersucht einen semi-supervisierten Ansatz zur automatischen Annotation von UAV-Videodaten mittels Transfer Learning, Pseudo-Labeling und tracking-basierter Label-Propagation. Als Datengrundlage dient der MBG-V2-Datensatz, der Drohnenaufnahmen urbaner Gebiete sowie Annotationen potenzieller Brutstätten enthält. Zunächst wird ein YOLOv8-basiertes Baseline-Modell auf manuell annotierten Frames trainiert. Anschließend werden zusätzliche Pseudo-Labels auf nicht manuell annotierten Frames erzeugt. Ergänzend wird ByteTrack eingesetzt, um Detektionen über mehrere aufeinanderfolgende Frames hinweg zeitlich zu verknüpfen und tracking-basierte Pseudo-Labels zu generieren.

Die experimentellen Ergebnisse zeigen, dass das ausschließlich auf manuellen Annotationen trainierte Baseline-Modell die beste Gesamtleistung erreicht. Auf dem gemeinsamen Validierungssplit erzielt es eine Precision von 0{,}831, einen Recall von 0{,}485, einen mAP@50-Wert von 0{,}572 und einen mAP@50--95-Wert von 0{,}343. Das detection-basierte Pseudo-Labeling erweitert zwar die Datenbasis und verbessert die Abdeckung mehrerer unterrepräsentierter Klassen, erreicht jedoch mit einer Precision von 0{,}692, einem Recall von 0{,}489, einem mAP@50-Wert von 0{,}564 und einem mAP@50--95-Wert von 0{,}323 keine klare Leistungssteigerung. Der leicht höhere Recall geht mit einem deutlichen Verlust an Precision einher.

Die tracking-basierte Label-Propagation führt zu konservativeren Pseudo-Labels und erzielt eine Precision von 0{,}765, einen Recall von 0{,}445, einen mAP@50-Wert von 0{,}543 und einen mAP@50--95-Wert von 0{,}327. Damit verbessert Tracking die Zuverlässigkeit der Pseudo-Label-Auswahl gegenüber dem rein detection-basierten Ansatz, übertrifft jedoch nicht das manuell trainierte Baseline-Modell. Insgesamt zeigen die Ergebnisse, dass semi-supervisierte Verfahren den manuellen Annotierungsaufwand reduzieren können, ihre Wirksamkeit jedoch stark von der Qualität der erzeugten Pseudo-Labels abhängt. Kleine Objektgrößen, komplexe Hintergründe, Klassenungleichgewicht und instabile Detektionen bleiben zentrale Herausforderungen.

\vfill

\noindent\textbf{Schlagwörter:} Objekterkennung, UAV-Videodaten, Drohnenaufnahmen, \textit{Aedes aegypti}, Mückenbrutstätten, YOLOv8, Transfer Learning, semi-supervised Learning, Pseudo-Labeling, Label-Propagation, Multi-Object-Tracking, ByteTrack

\newpage